\ifdefined\isincluded
\else
\documentclass[a4paper, 12pt]{report}
\usepackage[margin=0.8in]{geometry}
\usepackage{amsmath}
\usepackage{amsfonts}
\usepackage{amssymb}
\usepackage{graphicx}
\usepackage[english]{babel}
\usepackage[autostyle, english = american]{csquotes}
\MakeOuterQuote{"}
\usepackage{setspace}
\onehalfspacing
\usepackage{booktabs}
\usepackage{longtable}
\usepackage{tabularx}
\usepackage{array}
\usepackage{url}
\usepackage{cite}

\makeatletter
\let\oldincludegraphics\includegraphics
\renewcommand{\includegraphics}[2][]{%
  \IfFileExists{#2}{%
    \oldincludegraphics[#1]{#2}%
  }{%
    \IfFileExists{images/#2}{%
      \oldincludegraphics[#1]{images/#2}%
    }{%
      \fbox{\parbox[c][3cm][c]{0.8\linewidth}{\centering \texttt{[Image: #2]}}}%
    }%
  }%
}
\makeatother

\begin{document}
\chapter{System Architecture and Experimental Methodology}
\fi

\label{chap.3}
\section{Two-Tier System Topology and Agent Functional Roles}
I propose a two tier multi agent business system architecture which decomposes complex decision making into smaller chunks. It has three expert agents. The first tier is the Planner Agent, A1 . It figures out what the user really needs, checks operational constraints against company policy and builds a detailed execution plan.

The second level is execution and validation. The Planner continuously sends representations of plans to Tool-Caller Agent (A 2 ). A 2 maps each planned action to a particular API schema, resulting in structured JSON arguments. The Verifier Agent (A3) verifies the policy adherence and observes the external environment response. It also ensures that the changes in the database are in sync with the customer request and in case something goes wrong or does not add up, it sends back corrective feedback.

Figure 3.1 illustrates the overall system architecture, and compares the two operation approaches. First, the classic Multi-Round Text ReAct. The second is Recursive Latent ReAct, which routes continuous tensors among agents. For the text-based baseline, each handoff between agents maps to generating one natural language critique token at a time. This bloats the context and increases inference latency. Recursive Latent ReAct stores intermediate deliberation in continuous representation space, and only decodes tokens on the boundary of database execution.

\begin{figure}[h!]
\centering
\includegraphics[width=\linewidth]{System Architecture.png}
\caption{System architecture. Multi-Round Text ReAct with verbal critique loops vs Recursive Latent ReAct passing continuous tensors through neural adapters in a two-tier design.}
\label{fig:sys_arch}
\end{figure}

\section{Foundation Model Backbone, Local Serving Infrastructure, and Reproducibility}
All comparisons are on the same Qwen3.5-4B model [11] ($P = 4.4 \times 10^9$ total parameters), so the results reflect differences in architecture design, not base capacity. This model includes embedding layers with a vocabulary size of $248{,}320$, 32 hybrid Transformer layers interleaving Gated DeltaNet linear attention with gated full self-attention \cite{qwen2026team}, a hidden dimension of $d_h = 2560$, 16 attention heads, and an 8,000-token active context budget. The model base weights were quantized to 4-bit medium block precision using GGUF (Q4\_K\_M), which produced a tiny 2.74 GB VRAM footprint. During forward passes, 4-bit quantized weights are dequantized in order to utilize 16-bit floating-point (bfloat16) activation vectors and 32-bit floating-point (float32) accumulation in CUDA cuBLAS GEMM kernels. Crucially, all neural adapters ($R_{\text{in}}, R_{\text{out}}$), activation vectors ($h \in \mathbb{R}^{2560}$), layer normalizations, and gradient updates must function in unquantized floating-point precision (bfloat16/float32) to prevent quantization noise in latent routing.

To ensure precise reproducibility, we provide the hardware and software configurations. For tool wrapper unit testing and local prototyping, an NVIDIA GeForce RTX 3050 laptop GPU (4 GB GDDR6 VRAM, Ampere architecture, 60W TGP) was utilized. The entire benchmark evaluation matrix (912 execution trajectories) and offline neural adapter pretraining were performed on a dedicated compute node with an NVIDIA RTX A4000 GPU (16 GB GDDR6 VRAM) running Ubuntu 22.04 LTS with CUDA 12.4, Python 3.10, PyTorch 2.4.0, and \texttt{llama-cpp-python} (v0.2.89 with FlashAttention-2 and cuBLAS). By completely offloading all 32 layers to GPU memory (\texttt{n\_gpu\_layers=33}), host paging and network variability were eliminated.

Fully deterministic decoding for all agents using fixed random seed 42 and greedy decoding (temperature 0.0, top-p 1.0). Token budgets are strictly restricted to 100 tokens for Tool-Caller schema outputs and 120 tokens for Verifier critiques. Monotonic hardware timers are used to measure wall-clock latency at the microsecond level.

\section{Neural Recursive Adapters, Pretraining Protocol, and Tensor Routing}
Our latent architecture controls communication between agent roles via continuous hidden representations rather than discrete textual tokens. $h \in \mathbb{R}^{d_h}$ represents the hidden state vector that was obtained from the final Transformer layer (Layer 32), with the hidden dimension $d_h = 2560$. During recursive passes, decision-making is mediated by two specialized neural adapter modules that employ full precision floating-point arithmetic.

The InnerLink adapter, denoted by $R_{\text{in}}$, does intra-agent representation refinement by updating an agent's internal latent state over recursive steps while preserving its functional specialization. It is built as a two-layer residual Multi-Layer Perceptron with pre-layer and post-layer normalization:
\begin{equation}
    R_{\text{in}}(h) = \text{LayerNorm}_{\text{post}}\left( h + W_2 \, \sigma\left( W_1 \, \text{LayerNorm}_{\text{pre}}(h) + b_1 \right) + b_2 \right)
\end{equation}
where $W_1 \in \mathbb{R}^{2560 \times 2560}$ and $W_2 \in \mathbb{R}^{2560 \times 2560}$ are projection weight matrices, and $\sigma$ is the Gaussian Error Linear Unit (GELU) activation function. The residual skip connection ensures numerical stability over iterative recursive runs by preserving the original semantic representation.

The OuterLink adapter, denoted by $R_{\text{out}}$, manages cross-agent subspace projection. From the Planner to the Tool-Caller ($R_{\text{out}}^{\text{RT}}$), from the Tool-Caller to the Verifier ($R_{\text{out}}^{\text{TR}}$), and from the Verifier back to the Planner, it translates representations between various agent functional roles. Since different agent roles operate in distinct conceptual subspaces, $R_{\text{out}}$ combines a straight linear residual projection with a higher-dimensional non-linear transformation:
\begin{equation}
    R_{\text{out}}(h) = \text{LayerNorm}_{\text{tgt}}\left( W_{\text{res}} h + W_2 \, \sigma\left( W_1 \, \text{LayerNorm}_{\text{src}}(h) + b_1 \right) + b_2 \right)
\end{equation}
where $W_1 \in \mathbb{R}^{5120 \times 2560}$, $W_2 \in \mathbb{R}^{2560 \times 5120}$, $W_{\text{res}} \in \mathbb{R}^{2560 \times 2560}$, and $\sigma$ is GELU.

Adapter weight matrices were initialized via He normal initialization with zero-initialized bias vectors, and the direct residual projection matrix $W_{\text{res}}$ was initialized as the identity operator. To establish role-specific and cross-agent coordination prior to benchmark deployment, the adapters were pretrained on an offline deliberation trajectory dataset comprising 10,000 multi-agent reasoning demonstrations split into an 80\% training split (8,000 trajectories) and a 20\% validation split (2,000 trajectories) with a fixed random seed of 42. These pretraining demonstrations were created by teachers enforcing sequential planning, tool invocation, and reflection trails throughout generic multi-step problem solving. Crucially, the pretraining data did not contain or overlap with any of the 114 $\tau^2$-Bench retail customer support tasks, tool definitions, customer profiles, or target database states. This ensures that the pretraining corpus is kept separate from the evaluation environment to avoid benchmark contamination.

Using the AdamW optimizer ($\beta_1 = 0.9, \beta_2 = 0.95$, weight decay of zero), InnerLink training optimized role-specific latent feature extraction with an initial learning rate of $5 \times 10^{-4}$ decayed via a cosine schedule across twenty thousand steps with ten warmup steps and a batch size of two. The inner objective minimized a composite alignment loss:
\begin{equation}
    \mathcal{L}_{\text{inner}} = 1.0 \times \left(1 - \text{CosSim}(h_{\text{pred}}, h_{\text{target}})\right) + 0.1 \times \| h_{\text{pred}} - h_{\text{target}} \|_2^2
\end{equation}
where $h_{\text{target}}$ represents the role-specific activation vector extracted from particular planning or tool-calling demonstrator traces. OuterLink training jointly optimized cross-agent bridging across three recursive deliberation passes ($r=3$) with frozen inner adapters under the supervision of terminal cross-entropy token generation loss paired with intermediate critique loss ($\lambda = 0.1$) using AdamW (learning rate $5 \times 10^{-4}$, gradient clipping norm of one point zero). Model checkpoints were evaluated every 1,000 steps on the held-out validation split. The checkpoint that had the lowest validation composite loss---checkpoint-10000 for InnerLink and the lowest joint validation cross-entropy plus critique loss for OuterLink---was selected. Throughout the benchmark evaluation, all adapter settings were executed precisely when frozen in evaluation mode. The Planner generates a hidden tensor $h_1^{(r)}$, which is refined by $R_{\text{in}}$ and projected by $R_{\text{out}}^{\text{RT}}$ in order to condition the Tool-Caller without generating textual tokens. This entire process occurs in latent space. The Tool-Caller strictly decodes to structured JSON tokens in order to carry out the external API instruction. The execution observation is then sent to the Verifier via $R_{\text{out}}^{\text{TR}}$.

\section{Enterprise Benchmark Environment: $\tau^2$-Bench Retail Domain}
We evaluate autonomous decision-making in a stateful enterprise environment using the Retail domain of the $\tau^2$-Bench benchmark \cite{sierra2024tau}. Thirteen deterministic Python Application Programming Interface tools and an interactive SQLite relational database are used in the benchmark to replicate a production-level customer service environment.

The database consists of 500 customer profiles and 1,000 orders with multiple items. It also has large product catalogs with stock levels, warehouse locations and return tracking records. The executable tools include user identification, order look up, pending order updates, item returns, product exchanges, refund calculations and promotional discount validation. The environment is operated to strict company procedures. Sorry, we do not accept returns after thirty days, we do not change orders once they have been delivered, and there are restocking fees on certain categories of items. A formal definition of an enterprise workflow task is a 5-tuple:

The formal definition of an enterprise workflow task is a five-tuple:
\begin{equation}
    T = (u, p, s_0, \mathcal{A}, G)
\end{equation}
where $p$ denotes organizational business policies, $u$ is the user's natural language request, $s_0$ is the initial database state, $\mathcal{A}$ is the set of executable tools, and $G$ is the ground-truth terminal goal criteria.

\section{Operational Evaluation Metrics and Rubric}
To establish an objective, repeatable evaluation standard, execution rollouts are evaluated using a strict multi-component scoring rubric that is directly based on benchmark ground truth. Task evaluation generates several operational measures, including continuous scoring, binary objective attainment, clean execution dependability, and policy constraint adherence.

The continuous goal satisfaction score, denoted by $S_{\text{goal}} \in [0.0, 1.0]$, is a weighted combination of communicated factual accuracy and database action execution:
\begin{equation}
    S_{\text{goal}} = 0.6 \times S_{\text{info}} + 0.4 \times S_{\text{act}}
\end{equation}
where $S_{\text{info}}$ calculates the percentage of required ground-truth entities and status values that are correctly given in the final customer response, and $S_{\text{act}}$ displays the percentage of required state changes and database queries that are correctly executed. The task goal satisfaction accuracy score is represented as a percentage ($S_{\text{goal}} \times 100\% \in [0.0\%, 100.0\%]$). The per-depth and per-tier accuracy in Table 5.2 and Table 5.3 show the average accuracy of continuous goal satisfaction over all the tasks in each category. So they show fractional percentages instead of discrete jumps. The binary task success is recorded by the boolean flag success. It is rewarded if the score of goal satisfaction is greater than a threshold of 70\% ($S_{\text{goal}} \ge 0.70$). At terminal recursion level ($r=3$), 62 Text ReAct tasks (54.4\%) and 64 Latent ReAct tasks (56.1\%) out of 114 tasks achieve universal binary success. Clean success tasks are defined as the rigorous subset of tasks at terminal depth ($r=3$) that achieve binary success ($S_{\text{goal}} \ge 0.70$) without causing any tool execution exceptions, malformed argument structures, or policy violations during their rollout (30 tasks for Text ReAct versus 33 tasks for Latent ReAct). The arithmetic average of continuous goal fulfillment accuracy over all four recursion depths ($r=0\dots3$) is represented by the mean task completion accuracy across recursion depths (50.3\% vs. 50.9\%).
Computational inference expense in TeraFLOPs (TFLOPs) is analytically estimated using standard scaling formulae for Transformer autoregression \cite{vaswani2017attention}:
\begin{equation}
    \text{FLOPs}_{\text{task}} = \frac{2 \times P \times (N_{\text{prompt}} + N_{\text{gen}}) + \text{FLOPs}_{\text{adapters}}}{10^{12}}
\end{equation}
where $P = 4.4 \times 10^9$ represents the total model parameters (all parameters participate in each dense forward pass, as Qwen3.5-4B is a dense non-MoE architecture), $N_{\text{prompt}}$ is the total prompt tokens ingested across all roles (Planner, Tool-Caller, Verifier) and passes ($r \in \{0, 1, 2, 3\}$), and $N_{\text{gen}}$ is the total generated tokens. The neural adapters contribute about $2.6 \times 10^7$ $\text{FLOPs}(R_{\text{in}})$ and $6.5 \times 10^7$ $\text{FLOPs}(R_{\text{out}})$ FLOPs per pass ($< 0.001\%$ of backbone forward compute). The size of gated DeltaNet linear attention layers is $O(N \cdot d_h^2)$ without quadratic memory matrices in the hybrid architecture of Qwen3.5-4B \cite{qwen2026team}; for the subset of full self-attention layers inside our evaluated sequence budget ($N_{\text{turn}} \le 2048$ tokens), dense projections ($2P$ FLOPs/token) account for more than $97\%$ of theoretical computation. We thus ignore the minor quadratic attention term as a simple analytical approximation, recognizing this constraint while pointing out that precise runtime procedures rely on hardware caching. According to Table~\ref{tab:macro_performance}, cumulative analytical estimations for all 456 tasks per architecture add up to 13,076 TFLOPs for Text ReAct and 11,038 TFLOPs for Latent ReAct. precisely matching the total depth expenditures per run plotted in Figure~\ref{fig:flops_depth}.

Using the fixed benchmark denominator of 114 tasks (28 tasks or 24.6\%), benchmark error recovery yield is defined across the entire benchmark cohort ($N = 114$ tasks) as the percentage of all tasks that showed an initial-pass ($r=0$) execution failure or malformed invocation but successfully achieved binary goal completion ($\text{is\_success} = \text{True}$) upon subsequent recursive deliberation ($r \in \{1, 2, 3\}$; 24.6\% for Text ReAct as opposed to 33 assignments or 28.9\% for Latent ReAct). Policy violations are evaluated deterministically against business rules when an agent modifies immutable items (e.g., changing addresses on dispatched orders, giving excess refunds, or omitting identity verification). Isolated tool selection accuracy is the percentage of API invocations that correspond to valid ground-truth tool schemas.

\section{Composite Task Complexity Metric and Statistical Testing Methodology}
Enterprise client interactions range widely in complexity, from straightforward status inquiries to complicated multi-step transaction disputes. To give an impartial and systematic evaluation, we created a composite difficulty score metric $C$:
\begin{equation}
    C = 1.0 \times C_{\text{tool}} + 2.0 \times C_{\text{write}} + 1.0 \times C_{\text{div}}
\end{equation}
where the distinct tool diversity count is represented by $C_{\text{div}}$, the number of database state changes is represented by $C_{\text{write}}$, and the minimum required tool invocations is represented by $C_{\text{tool}}$.

Using the composite score $C$, the 114 tasks were split into three balanced tiers of 38 tasks each. The Easy Tier ($C \le 2$) includes read-only queries with one to two API calls. The Medium Tier ($3 \le C \le 5$) requires three to four API requests and one database modification (e.g., updating shipping details following user verification). The Hard Tier ($C \ge 6$) includes five or more API calls, several state changes, relational dependencies, and strict policy validations (such multi-item returns with restocking deductions and stock replacement checks).

In all systems, 912 full trajectories are generated by considering the stated recursion levels $r \in \{0, 1, 2, 3\}$. Formal inferential hypothesis testing is specifically concerned with representation dynamics: step-wise trajectory drift contraction and inter-agent directional alignment across the $N = 114$ task pairings are evaluated using paired, two-sided Wilcoxon signed-rank tests. Non-parametric effect sizes between successful and unsuccessful rollouts are measured using the Mann-Whitney $U$ test and Cliff's Delta ($\delta$). Family-wise error rates are tightly controlled across the family of $m = 10$ candidate geometric hypotheses using Holm-Bonferroni corrections. All indicated confidence intervals are non-parametric percentile intervals computed over 10,000 bootstrap resamples. However, total benchmark accuracy differences (+0.6 percentage points across depths, +2.0 percentage points at terminal depth) are reported as observed descriptive differences across the benchmark cohort instead of formal statistical significance claims.

\ifdefined\isincluded
\else
\begin{thebibliography}{99}
\bibitem{qwen2026team} Qwen Team, ``Qwen3.5: Hybrid Attention Language Models,'' \emph{Technical Report, Alibaba Group}, 2026.
\bibitem{sierra2024tau} S.~Sierra, M.~Ritter, K.~Narasimhan, and S.~Yao, ``$\tau$-Bench: A Benchmark for Tool-Agent-User Interaction in Real-World Domains,'' \emph{arXiv preprint arXiv:2406.12045}, 2024.
\bibitem{vaswani2017attention} A.~Vaswani, N.~Shazeer, N.~Parmar, J.~Uszkoreit, L.~Jones, A.~N. Gomez, {\L}.~Kaiser, and I.~Polosukhin, ``Attention is All You Need,'' in \emph{Advances in Neural Information Processing Systems (NeurIPS)}, 2017.
\end{thebibliography}
\end{document}
\fi
